\documentclass[12pt]{article}

\usepackage{sbc-template}
\usepackage{graphicx,url}
\usepackage{microtype}
\usepackage[utf8]{inputenc}
\usepackage[brazil]{babel}
\usepackage{comment}
\usepackage{lipsum}
\usepackage{csquotes}
\usepackage{url}
\usepackage[section]{placeins}
\usepackage[
    style=authoryear,
    citestyle=authoryear,
    backend=biber
]{biblatex}
\usepackage{pgfplots}

\pgfplotsset{compat=1.18}
\addbibresource{references.bib}

\DeclareFieldFormat{labelyear}{#1}
\renewcommand*{\nameyeardelim}{\addspace}
\renewbibmacro*{cite}{%
  \printtext[brackets]{%
    \printnames{labelname}%
    \setunit{\nameyeardelim}%
    \printfield{labelyear}%
  }}

\urlstyle{same} 

\sloppy

\title{Deteção de Deepfakes: Uma Abordagem com SigLIP2 e Explicabilidade via Grad-CAM}

\author{Jessica Fileto\inst{1}}

\address{
  Centro de Matemática, Computação e Cognição\\
  Universidade Federal do ABC (UFABC)\\
  Av. dos Estados, 5001 -- 09210-580 -- Santo André -- SP -- Brasil
  \email{jessica.fileto@ufabc.edu.br}
}

\begin{document}

\maketitle
     
\begin{resumo} 
Um dos grandes desafios atuais é a mitigação da desinformação, especialmente com o
advento de vídeos sintéticos hiper-realistas, conhecidos como deepfakes.
O desenvolvimento de técnicas de detecção de deepfakes é crucial para preservar a integridade da informação
e proteger indivíduos contra a disseminação de conteúdo falso.
Além disso, a explicabilidade dos modelos de detecção é essencial para aumentar a confiança dos usuários
e permitir uma compreensão mais profunda das decisões do modelo.
Este artigo propõe uma abordagem para a detecção de deepfakes utilizando SigLIP2
para extração de características visuais, juntamente com técnicas de explicabilidade,
como Grad-CAM, para fornecer insights sobre as decisões do modelo. A metodologia envolve o pré-processamento de dados,
extração de características com SigLIP2, seguida pela classificação utilizando redes neurais profundas.
A avaliação do desempenho do modelo é realizada através de métricas de precisão.
Para garantir a robustez do modelo, é realizado o teste da generalização em diferentes \textit{datasets},
sendo estes: Celeb-DF, DeeperForensics-1.0 e um \textit{dataset} customizado.
\end{resumo}

\begin{abstract}
One of the major challenges today is the mitigation of disinformation,
particularly with the advent of hyper-realistic synthetic videos known as deepfakes.
The development of deepfake detection techniques is crucial for preserving information
integrity and protecting individuals against the dissemination of false content.
Furthermore, the explainability of detection models is essential for increasing user trust and enabling a deeper understanding of model decisions.
This paper proposes an approach for deepfake detection using SigLIP2 for visual feature extraction,
together with explainability techniques such as Grad-CAM to provide insights into the model's decisions.
The methodology involves data preprocessing, feature extraction with SigLIP2, followed by classification using deep neural networks.
Model performance is evaluated through precision metrics.
To ensure model robustness, generalization testing is conducted across different datasets, namely: Celeb-DF, DeeperForensics-1.0, and a custom-built dataset.
\end{abstract}

\section{Introdução}

A divulgação de conteúdo é uma das formas mais comuns de comunicação,
principalmente com o advento das redes sociais. Mídias diversas podem ser distribuídas como foto, áudio e vídeo.

Com o acesso facilitado a ferramentas de inteligência artificial generativa (GenAI),
como por exemplo, Kling AI \cite{KlingAINextGeneration} , Runway \cite{RunwayBuildingAI} e o próprio Veo da Google \cite{Veo3Google},
os usuários conseguem através de imagens (IPV) e texto (TPV) gerar vídeos hiper realistas,
o que torna imperceptível ao ser humano detectar se foi gerado sinteticamente ou não.
Estes vídeos podem ser utilizados para espalhar desinformação e causar danos a sociedade como um todo.
Um dos métodos mais comuns de desinformação é o \textit{deepfake}, que consiste em vídeos ou imagens manipuladas com o intuito de enganar o público.
Os \textit{deepfakes} são criados utilizando \textit{deep learning}, que é uma técnica de aprendizado de máquina,
para substituir o rosto de uma pessoa por outro em um vídeo existente, 
ou criar um vídeo completamente novo com a aparência de uma pessoa real. \cite{dagarLiteratureReviewPerspectives2022a}

Apesar de passar credibilidade, as pessoas tendem a não acreditar na mensagem passada,
caso o conteúdo não seja condizente. Por exemplo, um político moderado ter falas de extrema direita,
porém caso a alteração seja sutil existe uma tendência a não ser percebido como \textit{deepfake}.
\cite{hameleersYouWontBelieve2022,hameleersTheyWouldNever2024}

Segundo \textcite[p.~10]{hameleersYouWontBelieve2022} “A desinformação que é congruente com as crenças anteriores das pessoas
pode reforçar as crenças existentes.“

A quantidade massiva de conteúdo sintético nas redes sociais ocasiona a fadiga de discernimento
da autenticidade do conteúdo que o usuário interage e impacta sua capacidade
e vontade de interagir criticamente com os meios de comunicação,
que propicia a aceitação passiva do conteúdo disponível. \cite{liuWhenSeeingNot2025a}

Segundo o relatório da \textit{Security Hero}, 98\% das \textit{deepfakes} online são de pornografia. \cite{2023StateDeepfakes}
Em relação ao Brasil, o relatório da Agência Lupa aborda o crescimento de 25,77\% 
no uso de GenIA na disseminação de desinformação no ano de 2025, sendo o tipo de conteúdo mais disseminado, o vídeo.
O maior alvo de desinformação continua sendo a polarização política, atingindo 50,2\% 
da quantidade de conteúdos nacionais analisados pela Agência Lupa em 2025. \cite{correaPanoramaDesinformacaoNo2026}
A \textit{SaferNet} fez um levantamento dos casos de deepfakes,
focados em deepfakes de pornografia em escolas brasileiras. Neste âmbito dos 27 estados foram constatados 10 estados
com relatos desta prática. \cite{MapeamentoSaferNetIdentificaa}

Os danos desta prática são imensos, visto que mesmo que ocorra a detecção do conteúdo \textit{fake},
no caso o deepfake de pornografia, isto não impede sua rápida disseminação e reenvio entre plataformas.
Por fim, saber que uma imagem é falsa não anula os danos ou perda de dignidade infligida à pessoa retratada.
\cite{mahrSexualizedDeepfakesSociotechnical2026}

\section{Trabalhos relacionados}
Com o aumento da criação de deepfakes, houve também uma crescente pesquisa relacionada aos métodos de detecção.
Inicialmente, principalmente no campo forense, eram utilizadas técnicas voltadas a análise pixel a pixel.
O objetivo é encontrar inconsistências frame a frame, relacionados ao piscar dos olhos,
postura, diferenças de tom de pele e inconsistências da iluminação. 

As técnicas tradicionais de aprendizado de máquina (ML), também são empregadas na detecção de deepfakes,
sendo as mais comumente utilizadas o perceptron multicamada (MLP), regressão logística (LR) dentre outros modelos que são o estado da arte. 

Com o avanço das técnicas de geração de deepfakes, foi necessário avançar nos métodos de detecção.
Os métodos mais utilizados atualmente são os modelos baseados em Deep Learning, como por exemplo, redes neurais convolucionais (CNN),
redes neurais recorrentes (RNN) e redes neurais convolucionais baseadas em região (RCNN). 

Conforme  “Após o sucesso dos transformadores em visão computacional (ViT), houve uma exploração do ViT e variantes espaço-temporais […]”.
Um desses modelos utilizados é o SigLIP , que por ser pré-treinado em um extenso dataset é costumeiramente utilizado como um backbone congelado para extração de características 

\section{Delimitação e Escopo}

\section{Metodologia}


\subsection{Divisão dos dados}


\subsection{Pré-Processamento dos dados}


\subsection{Análise do desempenho}


\section{Resultados e discussões}

\subsection{Métricas de performance}

\section{Conclusão}

\section{Trabalhos Futuros}

\printbibliography

\end{document}
